% GENERATED FILE. Edit canonical lessons, then run npm run generate:books.
% canonical-source-hash: fnv1a64:6bacbb7bae3abc34
% canonical-lessons: KA-C248-dasaka, KA-C248-varsikotsava, KA-C248-gaduvu, KA-C248-muhurta, KA-C248-onduvare

\chapter{Decade, Anniversary, Deadline, Auspicious time, One and a half}
\label{ch:ka-decade-anniversary-deadline-auspicious-time-one-and-a-half}
\hlchaptermodality{248}

\begin{chapteropening}
\textbf{By the end of this chapter:} \emph{I can say a decade, an anniversary, a deadline, an auspicious time and one and a half.}

Complete the last lesson of chapter 248: I can say a decade, an anniversary, a deadline, an auspicious time and one and a half.
\end{chapteropening}

\section[\texorpdfstring{\kn{ದಶಕ} (daśaka)}{ದಶಕ (daśaka)}]{\kn{ದಶಕ} (daśaka) — a decade}

\label{lesson:KA-C248-dasaka}



\begin{quote}
Before the new one: say the Kannada for Ganesha Chaturthi, then the Kannada for Rajyotsava.
\end{quote}

\subsection*{You'll want to know: \kn{ದಶಕ}}
\textbf{\kn{ದಶಕ}} — \emph{daśaka} — \textquotedblleft{}a decade\textquotedblright{}.

\begin{grammarlens}[title={What you've built}]
One more word for this chapter.
\end{grammarlens}

\subsection*{Your turn}
\begin{itemize}
\raggedright
  \item \emph{Say it:} \emph{daśaka}
  \item \emph{Say it:} \emph{daśaka}, once more
  \item \emph{Say it:} say \emph{daśaka}
  \item \emph{Recall:} say the Kannada for Ganesha Chaturthi, then the Kannada for Rajyotsava, then say \emph{daśaka} again
\end{itemize}

\begin{tcolorbox}[breakable,colback=teal!4,colframe=teal!35!black,title={Before you move on}]
What does \kn{ದಶಕ} mean? (\textquotedblleft{}A decade\textquotedblright{}.) Say it once more.
\end{tcolorbox}
\section[\texorpdfstring{\kn{ವಾರ್ಷಿಕೋತ್ಸವ} (vārṣikōtsava)}{ವಾರ್ಷಿಕೋತ್ಸವ (vārṣikōtsava)}]{\kn{ವಾರ್ಷಿಕೋತ್ಸವ} (vārṣikōtsava) — an anniversary; an annual day}

\label{lesson:KA-C248-varsikotsava}



\begin{quote}
Before the new one: say the Kannada for Rajyotsava, then the Kannada for a decade.
\end{quote}

\subsection*{You'll want to know: \kn{ವಾರ್ಷಿಕೋತ್ಸವ}}
\textbf{\kn{ವಾರ್ಷಿಕೋತ್ಸವ}} — \emph{vārṣikōtsava} — \textquotedblleft{}an anniversary; an annual day\textquotedblright{}.

\begin{grammarlens}[title={What you've built}]
One more word for this chapter.
\end{grammarlens}

\subsection*{Your turn}
\begin{itemize}
\raggedright
  \item \emph{Say it:} \emph{vārṣikōtsava}
  \item \emph{Say it:} \emph{vārṣikōtsava}, once more
  \item \emph{Say it:} say \emph{vārṣikōtsava}
  \item \emph{Recall:} say the Kannada for Rajyotsava, then the Kannada for a decade, then say \emph{vārṣikōtsava} again
\end{itemize}

\begin{tcolorbox}[breakable,colback=teal!4,colframe=teal!35!black,title={Before you move on}]
What does \kn{ವಾರ್ಷಿಕೋತ್ಸವ} mean? (\textquotedblleft{}An anniversary; an annual day\textquotedblright{}.) Say it once more.
\end{tcolorbox}
\section[\texorpdfstring{\kn{ಗಡುವು} (gaḍuvu)}{ಗಡುವು (gaḍuvu)}]{\kn{ಗಡುವು} (gaḍuvu) — a deadline}

\label{lesson:KA-C248-gaduvu}



\begin{quote}
Before the new one: say the Kannada for a decade, then the Kannada for an anniversary.
\end{quote}

\subsection*{You'll want to know: \kn{ಗಡುವು}}
\textbf{\kn{ಗಡುವು}} — \emph{gaḍuvu} — \textquotedblleft{}a deadline\textquotedblright{}.

\begin{grammarlens}[title={What you've built}]
One more word for this chapter.
\end{grammarlens}

\subsection*{Your turn}
\begin{itemize}
\raggedright
  \item \emph{Say it:} \emph{gaḍuvu}
  \item \emph{Say it:} \emph{gaḍuvu}, once more
  \item \emph{Say it:} say \emph{gaḍuvu}
  \item \emph{Recall:} say the Kannada for a decade, then the Kannada for an anniversary, then say \emph{gaḍuvu} again
\end{itemize}

\begin{tcolorbox}[breakable,colback=teal!4,colframe=teal!35!black,title={Before you move on}]
What does \kn{ಗಡುವು} mean? (\textquotedblleft{}A deadline\textquotedblright{}.) Say it once more.
\end{tcolorbox}
\section[\texorpdfstring{\kn{ಮುಹೂರ್ತ} (muhūrta)}{ಮುಹೂರ್ತ (muhūrta)}]{\kn{ಮುಹೂರ್ತ} (muhūrta) — an auspicious time}

\label{lesson:KA-C248-muhurta}



\begin{quote}
Before the new one: say the Kannada for an anniversary, then the Kannada for a deadline.
\end{quote}

\subsection*{You'll want to know: \kn{ಮುಹೂರ್ತ}}
\textbf{\kn{ಮುಹೂರ್ತ}} — \emph{muhūrta} — \textquotedblleft{}an auspicious time\textquotedblright{}.

\begin{grammarlens}[title={What you've built}]
One more word for this chapter.
\end{grammarlens}

\subsection*{Your turn}
\begin{itemize}
\raggedright
  \item \emph{Say it:} \emph{muhūrta}
  \item \emph{Say it:} \emph{muhūrta}, once more
  \item \emph{Say it:} say \emph{muhūrta}
  \item \emph{Recall:} say the Kannada for an anniversary, then the Kannada for a deadline, then say \emph{muhūrta} again
\end{itemize}

\begin{tcolorbox}[breakable,colback=teal!4,colframe=teal!35!black,title={Before you move on}]
What does \kn{ಮುಹೂರ್ತ} mean? (\textquotedblleft{}An auspicious time\textquotedblright{}.) Say it once more.
\end{tcolorbox}
\section[\texorpdfstring{\kn{ಒಂದೂವರೆ} (ondūvare)}{ಒಂದೂವರೆ (ondūvare)}]{\kn{ಒಂದೂವರೆ} (ondūvare) — one and a half}

\label{lesson:KA-C248-onduvare}



\begin{quote}
Before the new one: say the Kannada for a deadline, then the Kannada for an auspicious time.
\end{quote}

\subsection*{You'll want to know: \kn{ಒಂದೂವರೆ}}
\textbf{\kn{ಒಂದೂವರೆ}} — \emph{ondūvare} — \textquotedblleft{}one and a half\textquotedblright{}.

\begin{grammarlens}[title={What you've built}]
One more word for this chapter.
\end{grammarlens}

\subsection*{Your turn}
\begin{itemize}
\raggedright
  \item \emph{Say it:} \emph{ondūvare}
  \item \emph{Say it:} \emph{ondūvare}, once more
  \item \emph{Say it:} say \emph{ondūvare}
  \item \emph{Recall:} say the Kannada for a deadline, then the Kannada for an auspicious time, then say \emph{ondūvare} again
\end{itemize}

\begin{tcolorbox}[breakable,colback=teal!4,colframe=teal!35!black,title={Before you move on}]
What does \kn{ಒಂದೂವರೆ} mean? (\textquotedblleft{}One and a half\textquotedblright{}.) Say it once more.
\end{tcolorbox}
